\subsection{Plane Traces}

The trace center $\CenterTrace$ represents the point 
whose motion defines the translational component; for centroidal traces, 
$\CenterTrace = \CentroidLocation$.

\begin{proposition}[Polynomial structure of traced strains]
\label{prop:strain-polynomial}
For any trace operator $\mathcal{T}$ satisfying the plane section property \eqref{eq:trace-rigid-property}, 
the strains extracted from a Saint-Venant field are at most linear in $\reflenx$:
\begin{equation}
\bm{\gamma}(\reflenx) = \bm{\gamma}_0 + \reflenx \, \bm{\gamma}_1, 
\qquad 
\bm{\kappa}(\reflenx) = \bm{\kappa}_0 + \reflenx \, \bm{\kappa}_1
\label{eq:strain-polynomial}
\end{equation}
Moreover, the coefficients depend linearly on the Ieşan parameters and are 
independent of the rigid body modes:
\begin{equation}
\begin{pmatrix} \bm{\gamma}_0 \\ \bm{\kappa}_0 \\ \bm{\gamma}_1 \\ \bm{\kappa}_1 \end{pmatrix}
= \mathbf{E}_{\mathcal{T}} \, \IesanVector
\label{eq:strain-coefficients}
\end{equation}
for some matrix $\mathbf{E}_{\mathcal{T}} \in \mathbb{R}^{12 \times 6}$.
\end{proposition}

\begin{proof}
The Saint-Venant displacement \eqref{eq:iesan-general-solution} is polynomial 
in $\reflenx$ with degree at most three, where the cubic terms arise solely from 
the flexural contribution $-\tfrac{1}{6}\reflenx^3 \bm{b}_\Section$. 
Any trace satisfying \eqref{eq:trace-rigid-property} preserves the polynomial 
structure, yielding $\bm{u}(\reflenx)$ cubic and $\bm{\theta}(\reflenx)$ quadratic. 
Differentiation reduces the degree by one, giving the linear structure 
\eqref{eq:strain-polynomial}.

For the independence from rigid modes, observe that a rigid displacement 
$\hat{\bm{u}}_{\rm R} = \bm{\delta}_u + \bm{\delta}_\theta \times (\reflenx \FrameBasis + \bm{r})$ 
traces to
\begin{equation*}
\bm{u}_R(\reflenx) = \bm{\delta}_u + \reflenx (\bm{\delta}_\theta \times \FrameBasis) 
+ \bm{\delta}_\theta \times \CenterTrace, 
\qquad 
\bm{\theta}_R(\reflenx) = \bm{\delta}_\theta
\end{equation*}
The resulting strains are
\begin{equation*}
\bm{\gamma}_R = \bm{\delta}_\theta \times \FrameBasis + \FrameBasis \times \bm{\delta}_\theta 
= \mathbf{0}, 
\qquad 
\bm{\kappa}_R = \mathbf{0}
\end{equation*}
Thus the rigid modes contribute nothing to the extracted strains.
\end{proof}

\begin{proposition}[Universality of strain gradients]
\label{prop:universal-gradients}
The linear coefficients $(\bm{\gamma}_1, \bm{\kappa}_1)$ are independent of the 
choice of trace operator:
\begin{equation}
\bm{\gamma}_1 = -(\CentroidLocation \cdot \bm{b}_\Section) \FrameBasis, 
\qquad 
\bm{\kappa}_1 = -\FrameBasis \times \bm{b}_\Section
\label{eq:universal-gradients}
\end{equation}
The trace dependence resides entirely in the constant coefficients 
$(\bm{\gamma}_0, \bm{\kappa}_0)$.
\end{proposition}

\begin{proof}
From the Saint-Venant kinematics, the highest-order terms in $\reflenx$ are 
uniquely determined by the flexural loading parameter $\bm{b}_\Section$. 
Specifically, the $\reflenx^2$ contribution to $\bm{\theta}(\reflenx)$ is 
$-\tfrac{1}{2}\reflenx^2 (\FrameBasis \times \bm{b}_\Section)$ for any trace, 
being the only quadratic term in the Saint-Venant rotation field. 
Differentiation yields $\bm{\kappa}_1 = -\FrameBasis \times \bm{b}_\Section$. 
Similarly, the axial strain gradient follows from the flexure-induced 
axial displacement.
\end{proof}

\begin{Quote}
\begin{proposition}[Plane trace matrix and gauge transformations]
Let $\mathcal{S}_{\mathrm{SV}}$ denote the six-dimensional space of 
Saint-Venant particular solutions modulo rigid body motions, parameterized 
by the Ieşan parameter vector $\IesanVector \in \mathbb{R}^6$ as in 
\eqref{eq:def-iesan-params}. Let 
\[
\mathcal{E}_0 \cong \mathbb{R}^6
\]
denote the space of constant strain coefficients
\[
\bm{e}_0 \triangleq \begin{pmatrix} \bm{\gamma}_0 \\ \bm{\kappa}_0 \end{pmatrix}.
\]

Any trace operator $\mathcal{T}$ satisfying the plane section property 
\eqref{eq:trace-rigid-property} induces a unique linear map
\[
E_0 : \mathcal{S}_{\mathrm{SV}} \to \mathcal{E}_0,
\qquad
\hat{\bm{u}}_{\mathrm{P}}(\,\cdot\,;\IesanVector) \mapsto \bm{e}_0(\IesanVector),
\]
and with respect to the chosen bases of $\mathcal{S}_{\mathrm{SV}}$ and 
$\mathcal{E}_0$ this map is represented by a matrix 
$\mathbf{T}_{\mathcal{T}} \in \mathbb{R}^{6\times 6}$ defined by
\begin{equation}
\bm{e}_0 = \mathbf{T}_{\mathcal{T}}\,\IesanVector.
\label{eq:def-trace-matrix-general}
\end{equation}

If the restriction of $\mathcal{T}$ to the Saint-Venant class is 
non-degenerate in the sense that $\bm{e}_0(\IesanVector) = \mathbf{0}$ implies 
$\IesanVector = \mathbf{0}$, then $\mathbf{T}_{\mathcal{T}}$ is invertible.

Moreover, if we change coordinates in the parameter space and the strain 
space by invertible matrices
\[
\widetilde{\IesanVector} = \mathbf{P}\,\IesanVector,
\qquad
\widetilde{\bm{e}}_0 = \mathbf{Q}\,\bm{e}_0,
\]
then the trace matrix transforms as
\begin{equation}
\widetilde{\mathbf{T}}_{\mathcal{T}} 
= \mathbf{Q}\,\mathbf{T}_{\mathcal{T}}\,\mathbf{P}^{-1}.
\label{eq:trace-matrix-change-of-basis}
\end{equation}
In particular, for two traces $\mathcal{T}_1$ and $\mathcal{T}_2$ with 
invertible trace matrices $\TraceDeDpOne$ and $\mathbf{T}_2$ in the same 
parameter basis, the corresponding constant strain coefficients are related by
\begin{equation}
\bm{e}_0^{(2)} = \mathbf{G}_{12}\,\bm{e}_0^{(1)},
\qquad
\mathbf{G}_{12} \triangleq \mathbf{T}_2\,\TraceDeDpOne^{-1}.
\label{eq:gauge-transform-general}
\end{equation}
\end{proposition}

\end{Quote}

\subsubsection{The Trace Matrix}

In view of Propositions~\ref{prop:strain-polynomial} and \ref{prop:universal-gradients}, 
the gauge-dependent content of a trace is captured by a $6 \times 6$ matrix 
relating the constant strain coefficients to the Ieşan parameters.

\begin{definition}[Trace matrix]
The \emph{trace matrix} $\mathbf{T}_{\mathcal{T}} \in \mathbb{R}^{6 \times 6}$ 
associated with a trace operator $\mathcal{T}$ is defined by
\begin{equation}
\bm{e}_0 \triangleq \begin{pmatrix} \bm{\gamma}_0 \\ \bm{\kappa}_0 \end{pmatrix} 
= \mathbf{T}_{\mathcal{T}} \, \IesanVector
\label{eq:def-trace-matrix}
\end{equation}
\end{definition}

\begin{Old}
\begin{proposition}[Structure of the trace matrix]
\label{prop:trace-matrix-structure}
For any trace operator $\mathcal{T}$ satisfying the plane section property 
\eqref{eq:trace-rigid-property} with trace center at the centroid, the trace 
matrix has the block structure
\begin{equation}
\mathbf{T}_{\mathcal{T}} = \begin{pmatrix}
1 & \mathbf{0}^\top & 0 & \TraceNF^\top \\[4pt]
\mathbf{0} & \mathbf{0} & \mathbf{0} & \TraceFF \\[4pt]
0 & \mathbf{0}^\top & 1 & \TraceTF^\top \\[4pt]
\mathbf{0} & -\FrameBasis^{\times} & \mathbf{0} & \TraceMF
\end{pmatrix}
\label{eq:trace-matrix-structure}
\end{equation}
where the blocks have dimensions consistent with the decomposition 
$\IesanVector \in \mathbb{R}^{1+2+1+2}$ and $\bm{e}_0 \in \mathbb{R}^{1+2+1+2}$:
\begin{itemize}
\item $\TraceFF \in \mathbb{R}^{2 \times 2}$: the \emph{shear-flexure coupling tensor}, 
relating transverse shear strain $\bm{\gamma}_{\perp,0}$ to the flexure parameter $\bm{b}_{\Section}$
\item $\TraceMF \in \mathbb{R}^{2 \times 2}$: captures curvature modifications 
from Poisson and warping effects
\item $\TraceNF \in \mathbb{R}^{2}$: trace-dependent axial strain shift from flexure
\item $\TraceTF \in \mathbb{R}^{2}$: twist-flexure coupling arising from $c_\vartheta(\bm{b}_{\Section})$
\end{itemize}
The notation $-\FrameBasis^{\times}$ denotes the $2 \times 2$ restriction of 
the cross-product operator to the section plane, acting as 
$-\FrameBasis^{\times}\bm{a}_{\Section} = -\FrameBasis \times (\FrameBasis_\alpha a_{\Section,\alpha}) 
= \FrameBasis_\beta \epsilon_{\alpha\beta} a_{\Section,\alpha}$.
\end{proposition}

\begin{proof}
We establish the structure by examining how each Ieşan parameter contributes 
to the constant strain coefficients $\bm{e}_0 = (\epsilon_{0}, \bm{\gamma}_{\perp,0}, 
\varkappa_0, \bm{\kappa}_{\perp,0})$.

\paragraph{Row 1: Axial strain $\epsilon_{0}$.}
The axial component of the shear strain is $\epsilon = \bm{u}' \cdot \FrameBasis$. 
From the Saint-Venant kinematics, the axial displacement has the form
\begin{equation*}
\FrameBasis \cdot \hat{\bm{u}}(\reflenx, \bm{r}) = \reflenx \, a_\varepsilon 
+ \tfrac{1}{2}\reflenx^2 (\bm{r} - \CentroidLocation) \cdot \bm{b}_{\Section} 
+ \text{(warping terms independent of $\reflenx$)}
\end{equation*}
Any trace satisfying the plane section property extracts $\bm{u}(\reflenx)$ such that 
$\FrameBasis \cdot \bm{u}'(0) = a_\varepsilon + \bm{\eta}_\gamma \cdot \bm{b}_{\Section}$ 
where $\bm{\eta}_\gamma$ depends on how the trace weights the $\reflenx^2$ term 
(which vanishes at $\reflenx = 0$ but contributes through derivatives of warping). 
For centroidal traces, $\bm{\eta}_\gamma = \mathbf{0}$. 
Thus:
\begin{equation*}
\epsilon_{0} = a_\varepsilon + \bm{\eta}_\gamma \cdot \bm{b}_{\Section}
\end{equation*}
This establishes the first row: $(1, \mathbf{0}^\top, 0, \bm{\eta}_\gamma^\top)$.

\paragraph{Row 2: Transverse shear strain $\bm{\gamma}_{\perp,0}$.}
The transverse shear strain is $\bm{\gamma}_\perp = \mathbf{1}_{\Section}(\bm{u}' + \FrameBasis \times \bm{\theta})$. 
The only Ieşan parameter that produces nonzero transverse shear at $\reflenx = 0$ 
is $\bm{b}_{\Section}$, through the interaction of the flexural displacement with 
the traced rotation. Specifically:
\begin{itemize}
\item The parameter $a_\varepsilon$ produces purely axial displacement: no contribution.
\item The parameter $\bm{a}_{\Section}$ produces rotation $\bm{\theta} \sim -\reflenx(\FrameBasis \times \bm{a}_{\Section})$ 
and displacement $\bm{u} \sim -\tfrac{1}{2}\reflenx^2 \bm{a}_{\Section}$, giving 
$\bm{\gamma}_\perp = -\reflenx \bm{a}_{\Section} + \reflenx \bm{a}_{\Section} = \mathbf{0}$ 
(the Euler-Bernoulli cancellation).
\item The parameter $a_\vartheta + c_\vartheta$ produces torsional rotation with 
no transverse shear contribution.
\item The parameter $\bm{b}_{\Section}$ produces the flexural field with nonzero 
shear strain at $\reflenx = 0$ through warping contributions.
\end{itemize}
The precise form depends on the trace. Define $\TraceFF$ by 
$\bm{\gamma}_{\perp,0} = \TraceFF \bm{b}_{\Section}$. 
This establishes 
the second row: $(\mathbf{0}, \mathbf{0}, \mathbf{0}, \TraceFF)$.

\paragraph{Row 3: Twist $\varkappa_0$.}
The twist is $\varkappa = \bm{\theta}' \cdot \FrameBasis$. From the Saint-Venant solution:
\begin{itemize}
\item The parameter $a_\varepsilon$ produces no rotation: no contribution.
\item The parameter $\bm{a}_{\Section}$ produces rotation perpendicular to $\FrameBasis$: no contribution.
\item The parameter $a_\vartheta + c_\vartheta$ produces uniform twist: 
$\varkappa_0 = a_\vartheta + c_\vartheta$.
\item The parameter $\bm{b}_{\Section}$ affects twist through the flexure-induced 
torsion $c_\vartheta = c_\vartheta(\bm{b}_{\Section})$, which is linear in $\bm{b}_{\Section}$ 
by \eqref{eq:iesan-2-22}.
\end{itemize}
Define $\bm{\eta}_\vartheta \in \mathbb{R}^2$ by $c_\vartheta = \bm{\eta}_\vartheta \cdot \bm{b}_{\Section}$ 
(absorbing the dependence into the third Ieşan parameter gives the stated form). 
Thus: $\varkappa_0 = (a_\vartheta + c_\vartheta) + \bm{\eta}_\vartheta \cdot \bm{b}_{\Section}$, 
but since $c_\vartheta$ is already included in the third component of $\IesanVector$, 
we write $\varkappa_0 = 1 \cdot (a_\vartheta + c_\vartheta) + \bm{\eta}_\vartheta \cdot \bm{b}_{\Section}$ 
where $\bm{\eta}_\vartheta$ captures any additional trace-dependent twist-shear coupling. 
This establishes the third row: $(0, \mathbf{0}^\top, 1, \bm{\eta}_\vartheta^\top)$.

\paragraph{Row 4: Transverse curvature $\bm{\kappa}_{\perp,0}$.}
The transverse curvature is $\bm{\kappa}_\perp = \mathbf{1}_{\Section} \bm{\theta}'$. 
From the Saint-Venant solution:
\begin{itemize}
\item The parameter $a_\varepsilon$ produces no rotation: no contribution.
\item The parameter $\bm{a}_{\Section}$ produces $\bm{\theta} = -\reflenx(\FrameBasis \times \bm{a}_{\Section}) + \ldots$, 
giving $\bm{\kappa}_{\perp,0} = -\FrameBasis \times \bm{a}_{\Section} = -\FrameBasis^\times \bm{a}_{\Section}$.
\item The parameter $a_\vartheta + c_\vartheta$ produces axial rotation only: no contribution.
\item The parameter $\bm{b}_{\Section}$ produces rotation modifications through Poisson 
and warping effects, captured by $\TraceMF \bm{b}_{\Section}$.
\end{itemize}
This establishes the fourth row: $(\mathbf{0}, -\FrameBasis^\times, \mathbf{0}, \TraceMF)$.

\paragraph{Conclusion.}
Assembling the four rows yields the stated block structure \eqref{eq:trace-matrix-structure}.
\end{proof}
\end{Old}

\begin{remark}[Non-centroidal traces]
For a trace with center $\CenterTrace \neq \CentroidLocation$, 
additional couplings appear. In particular, the bending parameter $\bm{a}_{\Section}$ 
can contribute to $\epsilon_{0}$ through the term 
$(\CenterTrace - \CentroidLocation) \cdot \bm{a}_{\Section}$, 
modifying the $(1,2)$ block from $\mathbf{0}^\top$ to a nonzero $1 \times 2$ matrix.
\end{remark}

\begin{proposition}[Structure of the trace matrix: general case]
Let $\mathcal{T}$ be a trace operator satisfying the plane section property 
\eqref{eq:trace-rigid-property} with trace center $\CenterTrace \in \Section$. 
Let the coordinate origin be arbitrary, with centroid at $\CentroidLocation$, 
shear center at $\CenterShearLocation$, and twist center at $\CenterTwistLocation$. 
The trace matrix has the block structure
\begin{equation}
\mathbf{T}_{\mathcal{T}} = \begin{pmatrix}
1 & (\CenterTrace - \CentroidLocation)^\top & 0 & \TraceNF^\top \\[4pt]
\mathbf{0} & \TraceNM & \bm{\xi}_{\vartheta} & \TraceFF \\[4pt]
0 & \mathbf{0} %\TraceTM^\top 
& 1 & \TraceTF^\top \\[4pt]
\mathbf{0} & -\FrameBasis^{\times} %+ \TraceMM 
& \TraceMT & \TraceMF
\end{pmatrix}
% \label{eq:trace-matrix-general}
\end{equation}
where the blocks have dimensions consistent with 
$\IesanVector \in \mathbb{R}^{1+2+1+2}$ and $\bm{e}_0 \in \mathbb{R}^{1+2+1+2}$:

\begin{center}
\renewcommand{\arraystretch}{1.3}
\begin{tabular}{c|c|c|l}
Block & Size & Units & Physical origin \\ \hline
$(\bm{r} - \CentroidLocation)^\top$ & $1 \times 2$ & L &
Axial-bending coupling from non-centroidal trace \\
$\TraceNF$ & $2 \times 1$ & L$^{2}$ &
Axial strain shift from flexure \\
$\bm{\Xi}^{(a)}$ & $2 \times 2$ & &
Shear-bending coupling (vanishes for Euler-Bernoulli) \\
$\bm{\xi}_{\vartheta}$ & $2 \times 1$ & &
Shear-twist coupling from non-coincident twist/shear centers \\
$\TraceFF$ & $2 \times 2$ & L$^2$ &
Shear-flexure coupling (encoding the classical shear correction) \\
% $\TraceTM$ & $2 \times 1$ & &
% Twist-bending coupling from asymmetric sections \\
$\TraceTF$ & $2 \times 1$ & L &
Twist-flexure coupling from $c_\vartheta(\bm{b}_{\Section})$ \\
% $\TraceMM$ & $2 \times 2$ & &
% Curvature-bending modification from trace \\
$\TraceMT$ & $2 \times 1$ & &
Curvature-twist coupling \\
$\TraceMF$ & $2 \times 2$ & L &
Curvature-flexure coupling from Poisson/warping
\end{tabular}
\end{center}
\end{proposition}

\begin{proof}
We systematically determine each block by computing the constant strain 
coefficients $\bm{e}_0 = (\epsilon_{0}, \bm{\gamma}_{\perp,0}, \varkappa_0, \bm{\kappa}_{\perp,0})$ 
as functions of the Ieşan parameters $\IesanVector = (a_\varepsilon, \bm{a}_{\Section}, a_\vartheta + c_\vartheta, \bm{b}_{\Section})$.

\paragraph{Step 1: General form of traced fields.}
For a trace operator $\mathcal{T}$ with center $\CenterTrace$, 
the traced displacement and rotation extracted from a field $\hat{\bm{u}}$ are
\begin{equation}
\bm{u}(\reflenx) = \mathcal{T}[\hat{\bm{u}}]_{\rm trans} - \mathcal{T}[\hat{\bm{u}}]_{\rm rot} \times \CenterTrace, 
\qquad 
\bm{\theta}(\reflenx) = \mathcal{T}[\hat{\bm{u}}]_{\rm rot}
\label{eq:traced-fields-general}
\end{equation}
where the subtraction of $\bm{\theta} \times \CenterTrace$ ensures that 
$\bm{u}$ represents the displacement of the trace center.

\paragraph{Step 2: Contribution from $a_\varepsilon$ (extension).}
The extension mode from \eqref{eq:iesan-displacement-1} is
\begin{equation*}
\hat{\bm{u}}^{(\varepsilon)} = (\reflenx \FrameBasis - \nu \bm{r}) a_\varepsilon
\end{equation*}
This field is linear in $\bm{r}$ with no $\reflenx$-dependent rotation. 
For any trace satisfying the plane section property:
\begin{equation*}
\bm{u}^{(\varepsilon)}(\reflenx) = \reflenx \FrameBasis \, a_\varepsilon - \nu \CenterTrace a_\varepsilon, 
\qquad 
\bm{\theta}^{(\varepsilon)}(\reflenx) = \mathbf{0}
\end{equation*}
The Poisson contraction $-\nu \bm{r}$ is a rigid translation when averaged, 
contributing only to $\bm{u}(0)$, not to strains. Thus:
\begin{equation*}
\bm{\gamma}^{(\varepsilon)} = (\bm{u}^{(\varepsilon)})' + \FrameBasis \times \bm{\theta}^{(\varepsilon)} 
= \FrameBasis \, a_\varepsilon, 
\qquad 
\bm{\kappa}^{(\varepsilon)} = \mathbf{0}
\end{equation*}
This gives the first column of $\mathbf{T}_{\mathcal{T}}$: $(1, \mathbf{0}, 0, \mathbf{0})^\top$.

\paragraph{Step 3: Contribution from $\bm{a}_{\Section}$ (bending).}
The bending mode from \eqref{eq:iesan-displacement-1} is
\begin{equation*}
\hat{\bm{u}}^{(a)} = \Big( -\tfrac{1}{2}\reflenx^2 \mathbf{1} - \nu \operatorname{dev}(\bm{r} \otimes \bm{r}) 
+ \reflenx \FrameBasis \otimes \bm{r} \Big) \bm{a}_{\Section}
\end{equation*}
The rotation extracted by the trace is (using the plane section property on 
the $\reflenx$-linear term):
\begin{equation*}
\bm{\theta}^{(a)}(\reflenx) = -\reflenx (\FrameBasis \times \bm{a}_{\Section}) 
+ \bm{\theta}^{(a)}_{\rm warp}
\end{equation*}
where $\bm{\theta}^{(a)}_{\rm warp}$ arises from the trace applied to the 
Poisson term $-\nu \operatorname{dev}(\bm{r} \otimes \bm{r}) \bm{a}_{\Section}$, 
which is quadratic in $\bm{r}$ and hence not purely rigid.

The displacement, accounting for the trace center offset, is
\begin{equation*}
\bm{u}^{(a)}(\reflenx) = -\tfrac{1}{2}\reflenx^2 \bm{a}_{\Section} 
+ \reflenx (\CenterTrace \cdot \bm{a}_{\Section}) \FrameBasis 
+ \reflenx (\FrameBasis \times \bm{a}_{\Section}) \times \CenterTrace 
+ \bm{u}^{(a)}_{\rm Poisson}
\end{equation*}
where $\bm{u}^{(a)}_{\rm Poisson}$ is the trace of the Poisson term.

Computing the strains at $\reflenx = 0$:
\begin{align*}
\epsilon_{0}^{(a)} &= \FrameBasis \cdot (\bm{u}^{(a)})'|_{\reflenx=0} 
= (\CenterTrace \cdot \bm{a}_{\Section}) 
= (\CenterTrace - \CentroidLocation) \cdot \bm{a}_{\Section} + \CentroidLocation \cdot \bm{a}_{\Section}
\end{align*}
However, from the Ieşan resultant equations \eqref{eq:2-15}, the centroidal 
contribution $\CentroidLocation \cdot \bm{a}_{\Section}$ is absorbed into the 
relationship between $\bm{a}_{\Section}$ and the bending moment. 
The axial strain is:
\begin{equation*}
\epsilon_{0}^{(a)} = (\CenterTrace - \CentroidLocation) \cdot \bm{a}_{\Section}
\end{equation*}

For the transverse shear:
\begin{align*}
\bm{\gamma}_{\perp,0}^{(a)} &= \mathbf{1}_{\Section} \big[ (\bm{u}^{(a)})' + \FrameBasis \times \bm{\theta}^{(a)} \big]_{\reflenx=0} \\
&= \mathbf{1}_{\Section} \big[ -\reflenx \bm{a}_{\Section} + (\FrameBasis \times \bm{a}_{\Section}) \times \CenterTrace 
+ \reflenx (\FrameBasis \times \bm{a}_{\Section}) + \FrameBasis \times \bm{\theta}^{(a)}_{\rm warp} \big]_{\reflenx=0} \\
&= (\FrameBasis \times \bm{a}_{\Section}) \times \CenterTrace + \FrameBasis \times \bm{\theta}^{(a)}_{\rm warp}
\end{align*}
Using the vector identity $(\FrameBasis \times \bm{a}_{\Section}) \times \CenterTrace 
= (\FrameBasis \cdot \CenterTrace) \bm{a}_{\Section} - (\bm{a}_{\Section} \cdot \CenterTrace) \FrameBasis 
= -(\bm{a}_{\Section} \cdot \CenterTrace) \FrameBasis$, the first term is axial only. 
Thus:
\begin{equation*}
\bm{\gamma}_{\perp,0}^{(a)} = \TraceNM \bm{a}_{\Section}
\end{equation*}
where $\TraceNM$ encodes the trace-dependent contribution 
from $\bm{\theta}^{(a)}_{\rm warp}$. 


For the curvatures:
\begin{equation*}
\varkappa_0^{(a)} = \FrameBasis \cdot (\bm{\theta}^{(a)})'|_{\reflenx=0} = 0 + \bm{\eta}_\kappa \cdot \bm{a}_{\Section}
\end{equation*}
where $\bm{\eta}_\kappa$ arises from any $\reflenx$-dependence in $\bm{\theta}^{(a)}_{\rm warp}$ 
(typically zero for standard traces).

\begin{equation*}
\bm{\kappa}_{\perp,0}^{(a)} = \mathbf{1}_{\Section} (\bm{\theta}^{(a)})'|_{\reflenx=0} 
= -\FrameBasis \times \bm{a}_{\Section} + \TraceMM \bm{a}_{\Section}
\end{equation*}
where $\TraceMM$ captures trace-dependent modifications.

This establishes the second column: 
$((\CenterTrace - \CentroidLocation)^\top, \TraceNM, \bm{\eta}_\kappa^\top, -\FrameBasis^\times + \TraceMM)^\top$.

\paragraph{Step 4: Contribution from $a_\vartheta + c_\vartheta$ (torsion).}
The torsion mode from \eqref{eq:iesan-displacement-1} is
\begin{equation*}
\hat{\bm{u}}^{(\vartheta)} = \big( \reflenx (\FrameBasis \times \bm{r}) + \WarpTwistIesan \FrameBasis \big) (a_\vartheta + c_\vartheta)
\end{equation*}
The traced fields are:
\begin{align*}
\bm{\theta}^{(\vartheta)}(\reflenx) &= \reflenx \FrameBasis (a_\vartheta + c_\vartheta) 
+ \RotationAverage{\WarpTwistIesan \FrameBasis} (a_\vartheta + c_\vartheta) \\
&= \big( \reflenx \FrameBasis + \CenterTwistLocation \big) (a_\vartheta + c_\vartheta)
\end{align*}
where $\CenterTwistLocation = \RotationAverage{\WarpTwistIesan \FrameBasis}$ is the 
twist center (a vector in the section plane).

\begin{align*}
\bm{u}^{(\vartheta)}(\reflenx) &= \reflenx (\FrameBasis \times \CenterTrace) (a_\vartheta + c_\vartheta) 
+ \SectionAverage{\WarpTwistIesan} \FrameBasis (a_\vartheta + c_\vartheta)  - \big( \reflenx \FrameBasis + \CenterTwistLocation \big) \times \CenterTrace (a_\vartheta + c_\vartheta) \\
&= \reflenx (\FrameBasis \times \CenterTrace) (a_\vartheta + c_\vartheta) 
- \CenterTwistLocation \times \CenterTrace (a_\vartheta + c_\vartheta)
\end{align*}
using $\SectionAverage{\WarpTwistIesan} = 0$ (Saint-Venant normalization) 
and $\FrameBasis \times \CenterTrace = -\CenterTrace \times \FrameBasis$.

Computing strains:
\begin{align*}
\epsilon_{0}^{(\vartheta)} &= 0 \\
\bm{\gamma}_{\perp,0}^{(\vartheta)} &= \mathbf{1}_{\Section} \big[ (\FrameBasis \times \CenterTrace) 
+ \FrameBasis \times (\FrameBasis + \tfrac{1}{\reflenx}\CenterTwistLocation) \big]_{\reflenx \to 0} (a_\vartheta + c_\vartheta)
\end{align*}
The term $\FrameBasis \times \FrameBasis = \mathbf{0}$, and 
$\FrameBasis \times \CenterTrace + \FrameBasis \times \CenterTwistLocation/\reflenx$ 
requires careful analysis. At $\reflenx = 0$:
\begin{equation*}
\bm{\gamma}_{\perp,0}^{(\vartheta)} = \FrameBasis \times (\CenterTrace - \CenterTwistLocation) (a_\vartheta + c_\vartheta) 
\triangleq \bm{\xi}_\vartheta (a_\vartheta + c_\vartheta)
\end{equation*}
where $\bm{\xi}_\vartheta = \FrameBasis \times (\CenterTrace - \CenterTwistLocation) \in \mathbb{R}^2$.

This coupling vanishes when $\CenterTrace = \CenterTwistLocation$ 
(trace center coincides with twist center).

For curvatures:
\begin{align*}
\varkappa_0^{(\vartheta)} &= \FrameBasis \cdot (\bm{\theta}^{(\vartheta)})' = a_\vartheta + c_\vartheta \\
\bm{\kappa}_{\perp,0}^{(\vartheta)} &= \mathbf{1}_{\Section} (\bm{\theta}^{(\vartheta)})' 
= \bm{\xi}_\vartheta^{(\kappa)} (a_\vartheta + c_\vartheta)
\end{align*}
where $\bm{\xi}_\vartheta^{(\kappa)}$ arises from any transverse component of 
$(\CenterTwistLocation)' = \mathbf{0}$, so typically $\bm{\xi}_\vartheta^{(\kappa)} = \mathbf{0}$.

This establishes the third column: $(0, \bm{\xi}_\vartheta, 1, \bm{\xi}_\vartheta^{(\kappa)})^\top$.

\paragraph{Step 5: Contribution from $\bm{b}_{\Section}$ (flexure).}
The flexure mode from \eqref{eq:iesan-displacement-2} is
\begin{align*}
\hat{\bm{u}}^{(b)} &= \Big( -\tfrac{1}{6}\reflenx^3 \mathbf{1}_{\Section} 
+ \reflenx \bm{\Pi}_{(b)} 
+ \tfrac{1}{2}\reflenx^2 \FrameBasis \otimes (\bm{r} - \CentroidLocation) 
+ \FrameBasis \otimes \bm{\WarpShearIesan} \Big) \bm{b}_{\Section} \\
&\quad + \big( \reflenx (\FrameBasis \times \bm{r}) + \WarpTwistIesan \FrameBasis \big) c_\vartheta(\bm{b}_{\Section})
\end{align*}

This is the most complex mode. The traced rotation is:
\begin{align*}
\bm{\theta}^{(b)}(\reflenx) &= -\tfrac{1}{2}\reflenx^2 (\FrameBasis \times \bm{b}_{\Section}) 
+ \reflenx \RotationAverage{\bm{\Pi}_{(b)} \bm{b}_{\Section}} 
+ \RotationAverage{\FrameBasis \otimes \bm{\WarpShearIesan}} \bm{b}_{\Section} + \big( \reflenx \FrameBasis + \CenterTwistLocation \big) c_\vartheta(\bm{b}_{\Section})
\end{align*}

The traced displacement (after subtracting $\bm{\theta}^{(b)} \times \CenterTrace$) is:
\begin{align*}
\bm{u}^{(b)}(\reflenx) &= -\tfrac{1}{6}\reflenx^3 \bm{b}_{\Section} 
+ \reflenx \SectionAverage{\bm{\Pi}_{(b)}} \bm{b}_{\Section} 
+ \tfrac{1}{2}\reflenx^2 (\CenterTrace - \CentroidLocation) \cdot \bm{b}_{\Section} \FrameBasis \\
&\quad + \text{(terms from $\bm{\theta}^{(b)} \times \CenterTrace$)} 
+ \text{($c_\vartheta$ contributions)}
\end{align*}

Computing strains at $\reflenx = 0$:

\textit{Axial strain:}
\begin{equation*}
\epsilon_{0}^{(b)} = \bm{\eta}_\gamma \cdot \bm{b}_{\Section}
\end{equation*}
where $\bm{\eta}_\gamma$ collects contributions from evaluating $(\bm{u}^{(b)})'$ 
at $\reflenx = 0$, including trace-dependent warping effects.

\textit{Transverse shear:}
\begin{equation*}
\bm{\gamma}_{\perp,0}^{(b)} = \TraceFF \bm{b}_{\Section}
\end{equation*}
where the shear-flexure tensor is
\begin{equation}
\TraceFF = \SectionAverage{\bm{\Pi}_{(b)}} 
+ \FrameBasis \times \RotationAverage{\FrameBasis \otimes \bm{\WarpShearIesan}} 
+ \FrameBasis \times (\CenterTrace - \CenterTwistLocation) \otimes \bm{\eta}_\vartheta 
+ \text{(trace-dependent terms)}
\label{eq:xi-general}
\end{equation}
The third term arises from the $c_\vartheta$ contribution to shear.

\textit{Twist:}
\begin{equation*}
\varkappa_0^{(b)} = \bm{\eta}_\vartheta \cdot \bm{b}_{\Section}
\end{equation*}
where, from \eqref{eq:iesan-2-22}:
\begin{equation}
\bm{\eta}_\vartheta = -\frac{1}{\Jo} \FrameBasis \cdot \int_{\Section} \bm{r} \times 
\big( \nabla\bm{\WarpShearIesan} + \bm{\Pi}_{(b)} \big) \mathrm{d}A
\label{eq:eta-vartheta-general}
\end{equation}

\textit{Transverse curvature:}
\begin{equation*}
\bm{\kappa}_{\perp,0}^{(b)} = \TraceMF \bm{b}_{\Section}
\end{equation*}
where $\TraceMF$ collects the $\reflenx$-independent 
contributions to $(\bm{\theta}^{(b)})'$.

This establishes the fourth column: 
$(\bm{\eta}_\gamma^\top, \TraceFF, \bm{\eta}_\vartheta^\top, \TraceMF)^\top$.

\paragraph{Step 6: Assembly.}
Collecting the four columns yields the stated block structure \eqref{eq:trace-matrix-general}.
\end{proof}

\begin{corollary}[Simplifications under geometric coincidences]
\label{cor:simplifications}
The trace matrix simplifies under the following geometric conditions:
\begin{enumerate}
\item \textbf{Centroidal trace} ($\CenterTrace = \CentroidLocation$): 
The $(1,2)$ block vanishes: $(\CenterTrace - \CentroidLocation)^\top = \mathbf{0}^\top$.

\item \textbf{Trace at twist center} ($\CenterTrace = \CenterTwistLocation$): 
The shear-twist coupling vanishes: $\bm{\xi}_\vartheta = \mathbf{0}$.

\item \textbf{Doubly-symmetric section}: 
The twist-flexure coupling vanishes: $\bm{\eta}_\vartheta = \mathbf{0}$, 
and $\TraceNM = \TraceMM = \mathbf{0}$.

\item \textbf{Pointwise trace at centroid} (Euler-Bernoulli): 
Additionally $\TraceFF = \mathbf{0}$, giving a block-diagonal structure.

\item \textbf{All centers coincident} 
($\CentroidLocation = \CenterShearLocation = \CenterTwistLocation = \CenterTrace$): 
The matrix reduces to
\begin{equation}
\mathbf{T}_{\mathcal{T}} = \begin{pmatrix}
1 & \mathbf{0} & 0 & \TraceNF^\top \\
\mathbf{0} & \mathbf{0} & \mathbf{0} & \TraceFF \\
0 & \mathbf{0} & 1 & \mathbf{0} \\
\mathbf{0} & -\FrameBasis^{\times} & \mathbf{0} & \TraceMF
\end{pmatrix}
\label{eq:trace-matrix-symmetric}
\end{equation}
\end{enumerate}
\end{corollary}

\begin{remark}[Physical interpretation of off-diagonal blocks]
The off-diagonal blocks in \eqref{eq:trace-matrix-general} encode the 
various couplings that arise in asymmetric sections:
\begin{itemize}
\item $(\CenterTrace - \CentroidLocation)^\top$: Bending of a non-centroidal 
fiber produces axial strain (the parallel axis effect).

\item $\TraceNM$: Pure bending produces apparent shear strain 
when the trace does not coincide with the neutral axis of warping.

\item $\bm{\xi}_\vartheta$: Pure twist produces apparent shear when the trace 
center differs from the twist center.

\item $\bm{\eta}_\kappa$: Bending produces apparent twist for sections where 
the principal bending and torsion axes are coupled.

\item $\bm{\eta}_\vartheta$: Flexure produces twist through the Ieşan coupling 
$c_\vartheta(\bm{b}_{\Section})$, present whenever the shear center does not 
coincide with the centroid.
\end{itemize}
These couplings are not artifacts of the beam theory but reflect genuine 
features of the three-dimensional Saint-Venant solution.
\end{remark}

\begin{remark}[Invertibility]
The trace matrix $\mathbf{T}_{\mathcal{T}}$ is invertible provided that 
$\TraceFF$ is invertible. 
This fails for the pointwise trace, 
where $\TraceFF = \mathbf{0}$, reflecting the fact that 
Euler-Bernoulli theory cannot independently prescribe shear strain. 
For shear-deformable theories (Timoshenko and its variants), 
$\TraceFF$ is generically invertible.
\end{remark}

The gauge parameters $(\bm{\eta}_\gamma, \eta_\vartheta, \TraceFF)$ 
of \S\ref{sec:sv} are now recognized as components of the trace matrix. 
Different trace choices correspond to different specifications of these components.

\subsubsection{Gauge Transformations}

The transformation between traces admits a simple algebraic characterization.

\begin{proposition}[Gauge transformation]
\label{prop:gauge-transform}
Let $\mathcal{T}_1$ and $\mathcal{T}_2$ be two trace operators with associated 
matrices $\TraceDeDpOne$ and $\mathbf{T}_2$. The constant strain coefficients 
transform as
\begin{equation}
\bm{e}_0^{(2)} = \mathbf{T}_2 \TraceDeDpOne^{-1} \, \bm{e}_0^{(1)} 
\triangleq \mathbf{G}_{12} \, \bm{e}_0^{(1)}
\label{eq:gauge-transform}
\end{equation}
where $\mathbf{G}_{12} = \mathbf{T}_2 \TraceDeDpOne^{-1}$ is the \emph{gauge transformation matrix}.
\end{proposition}

\begin{proof}
From $\bm{e}_0^{(1)} = \TraceDeDpOne \IesanVector$ and $\bm{e}_0^{(2)} = \mathbf{T}_2 \IesanVector$, 
eliminate $\IesanVector = \TraceDeDpOne^{-1} \bm{e}_0^{(1)}$.
\end{proof}

\subsubsection{The Section Constitutive Law}

The trace matrix provides a direct route to the section constitutive law.

\begin{proposition}[Trace-dependent constitutive law]
\label{thm:constitutive-law}
For a trace operator $\mathcal{T}$ with \emph{nonsingular} trace matrix $\mathbf{T}_{\mathcal{T}}$, 
the section constitutive law relating resultants to constant strain coefficients is
\begin{equation}
\begin{pmatrix} N \\ T \\ \bm{F} \\ \FrameBasis \times \bm{M} \end{pmatrix}
= \mathbf{C}_{\mathcal{T}} \, \bm{e}_0, 
\qquad 
\mathbf{C}_{\mathcal{T}} = \mathbf{R} \, \mathbf{T}_{\mathcal{T}}^{-1}
\label{eq:constitutive-law}
\end{equation}
where $\mathbf{R}$ is the resultant matrix \eqref{eq:def-R-matrix}.
\end{proposition}

\begin{proof}
Combining \eqref{eq:resultant-map} and \eqref{eq:def-trace-matrix}:
\begin{equation*}
\bm{s} = \mathbf{R} \IesanVector = \mathbf{R} \mathbf{T}_{\mathcal{T}}^{-1} \bm{e}_0
\end{equation*}
\end{proof}

\begin{corollary}[Transformation of constitutive matrices]
\label{cor:constitutive-transform}
The constitutive matrices for two traces are related by
\begin{equation}
\mathbf{C}_{\mathcal{T}_2} = \mathbf{C}_{\mathcal{T}_1} \, \mathbf{G}_{12}^{-1} 
= \mathbf{C}_{\mathcal{T}_1} \, \TraceDeDpOne \, \mathbf{T}_2^{-1}
\label{eq:constitutive-transform}
\end{equation}
\end{corollary}

This formula makes explicit the relationship between different beam theories: 
they share the same underlying Saint-Venant solution but extract different 
strain measures, leading to constitutive matrices related by a change of variables.

\begin{remark}[Symmetry of the constitutive matrix]
The constitutive matrix $\mathbf{C}_{\mathcal{T}}$ is symmetric if and only if 
$\mathbf{R} \mathbf{T}_{\mathcal{T}}^{-1}$ is symmetric. Since $\mathbf{R}$ 
encodes the Saint-Venant energy, which is symmetric, the symmetry of 
$\mathbf{C}_{\mathcal{T}}$ depends on the trace choice:
\begin{itemize}
\item The energy-conjugate trace $\mathcal{T}^{\rm conj}$ yields a symmetric 
$\mathbf{C}$ by construction.
\item The section-average trace $\mathcal{T}^{\rm avg}$ yields a symmetric 
$\mathbf{C}$ only for sections with sufficient symmetry.
\item The pointwise trace $\mathcal{T}^{\rm pt}$ typically yields a symmetric 
$\mathbf{C}$ but with singular shear blocks.
\end{itemize}
\end{remark}



\subsection{Summary}

The trace operator framework provides a unified view of shear-deformable beam theories:

\begin{enumerate}
\item A trace operator $\mathcal{T}$ extracts one-dimensional kinematic fields 
$(\bm{u}, \bm{\theta})$ from the three-dimensional field. 
The resulting strains  are linear in $\reflenx$, with universal gradients and trace-dependent constants.

\item The trace matrix $\mathbf{T}_{\mathcal{T}} \in \mathbb{R}^{6\times 6}$ 
encodes the relationship $\bm{e}_0 = \mathbf{T}_{\mathcal{T}} \IesanVector$.

\item The section constitutive law is $\mathbf{C}_{\mathcal{T}} = \mathbf{R} \mathbf{T}_{\mathcal{T}}^{-1}$.

\item Different traces yield constitutive matrices related by 
$\mathbf{C}_2 = \mathbf{C}_1 \TraceDeDpOne \mathbf{T}_2^{-1}$.

\item The trace choice determines which three-dimensional boundary conditions 
are representable by the one-dimensional model.
\end{enumerate}
